% ======================================================================
% Notation used by the game pseudocode of Lemma~\ref{lem:fs-hints}.
%
% This file is shared by two consumers:
%   * the paper, which \input{}s it from extended_proof.tex, and
%   * the TeXFrog HTML viewer, which \input{}s it via \tfmacrofile.
%
% Everything is \providecommand so that the paper's own macros.tex wins
% when both are loaded; the definitions below are the fallbacks that make
% the HTML build self-contained.
% ======================================================================

% The HTML viewer feeds this file to MathJax as well, to typeset the game
% names and descriptions in its sidebar.  MathJax knows neither \ensuremath
% nor \xspace, both of which the definitions below use; these two lines give
% it harmless equivalents.  Under real LaTeX both are already defined, so
% \providecommand does nothing.
\providecommand{\ensuremath}[1]{#1}
\providecommand{\xspace}{}

% --- Shared with report/macros.tex (fallbacks for the HTML build) ------
\providecommand{\secp}{\ensuremath{\lambda}}
\providecommand{\msg}{\textsf{msg}}
\providecommand{\com}{\mathsf{com}}
\providecommand{\chall}{\mathsf{chl}}
\providecommand{\rsp}{\mathsf{rsp}}
\providecommand{\hint}{\ensuremath{h}}
\providecommand{\hintdistr}{\ensuremath{\mathcal{H}}}
\providecommand{\numRO}{\ensuremath{q}}
\providecommand{\numOSign}{\ensuremath{s}}
\providecommand{\btrue}{{\texttt{true}}\xspace}
\providecommand{\bfalse}{\texttt{false}\xspace}
\providecommand{\until}{, \ldots,}
\providecommand{\cC}{\mathcal{C}}
\providecommand{\cE}{\mathcal{E}}
\providecommand{\Gen}{\mathsf{Gen}}
\providecommand{\Sim}{\mathcal{S}}
\providecommand{\Ver}{\mathsf{Ver}}
\providecommand{\bG}{\textbf{G}}
\providecommand{\responseupperbound}{\ensuremath{D_{\rsp}}}
\providecommand{\sqisign}{\textsc{SQIsign}\xspace}

% --- Specific to this section -----------------------------------------
\providecommand{\RO}{\ensuremath{\mathsf{RO}}}
\providecommand{\OSign}{\ensuremath{\mathsf{OSign}}}
\providecommand{\OTrans}{\ensuremath{\mathsf{OTrans}}}
\providecommand{\HT}{\ensuremath{\mathsf{HT}}}
\providecommand{\QT}{\ensuremath{\mathsf{QT}}}
\providecommand{\SIG}{\ensuremath{\mathsf{SIG}}}
\providecommand{\Prover}{\ensuremath{\mathcal{P}}}
\providecommand{\Verif}{\ensuremath{\mathcal{V}}}
\providecommand{\Real}{\ensuremath{\mathsf{Real}}}
\providecommand{\HintSim}{\ensuremath{\mathsf{HintSim}}}
\providecommand{\badev}{\ensuremath{\mathsf{bad}}}
\providecommand{\badprog}{\ensuremath{\mathsf{bad}_\mathsf{prog}}}
\providecommand{\goodev}{\ensuremath{\mathsf{Good}}}
\providecommand{\fpguess}{\ensuremath{\mathsf{fp}}}
\providecommand{\hctr}{\ensuremath{\mathsf{hc}}}
\providecommand{\sctr}{\ensuremath{\mathsf{sc}}}
\providecommand{\Qset}{\ensuremath{\mathcal{Q}}}
\providecommand{\hintimppaadv}[4]{%
  \ensuremath{\Advantage{\mathsf{hint}\text{-}\mathsf{IMP}\text{-}\mathsf{PA}}{#1,#2}(#3,#4)}}
\providecommand{\Advantage}[2]{\mathsf{Adv}^{#1}_{#2}}
\providecommand{\EUFCMA}{\ensuremath{\mathsf{EUF}\text{-}\mathsf{CMA}}\xspace}

% Explicit, game-independent line numbers, typeset exactly like the ones
% cryptocode's `linenumbering` option produces.  They are written by hand
% rather than generated because the proof cites individual lines and those
% citations have to denote the same line in every game of the chain.
\providecommand{\lno}[1]{\pclnstyle{#1}\pclnseparator\hspace{\pclnspace}}

% --- Macros cryptocode itself claims -----------------------------------
% The TeXFrog HTML wrapper loads cryptocode with the option set
%   [n,advantage,operators,sets,adversary,landau,probability,notions,logic,
%    ff,mm,primitives,events,complexity,oracles,asymptotics,keys]
% which defines \sig, \pk, \adv, \bdv, \ddv, \set, \concat, \sample, \abs and
% \minentropy itself.  A plain \providecommand would therefore be a no-op in
% the HTML build and the browser would disagree with the paper (cryptocode's
% \sig, for instance, is \mathsf{Sig}, not the signature variable we want).
% These ten are consequently defined unconditionally.  The values are exactly
% those of report/macros.tex, so re-running this inside the paper changes
% nothing there.
\providecommand{\pk}{}\renewcommand{\pk}{\textsf{pk}}
\providecommand{\adv}{}\renewcommand{\adv}{\mathcal{A}}
\providecommand{\bdv}{}\renewcommand{\bdv}{\mathcal{B}}
\providecommand{\ddv}{}\renewcommand{\ddv}{\mathcal{D}}
\providecommand{\sig}{}\renewcommand{\sig}{\ensuremath{\mathsf{sig}}}
\providecommand{\concat}{}\renewcommand{\concat}{\mathrel{||}}
\providecommand{\sample}{}\renewcommand{\sample}{%
  \gets\mathllap{\vcenter{\hbox{\tiny\rmfamily\upshape\$}}}\;}
\providecommand{\set}[1]{}\renewcommand{\set}[1]{\left\{#1\right\}}
\providecommand{\abs}[1]{}\renewcommand{\abs}[1]{\ensuremath{\lvert #1 \rvert}}
\providecommand{\minentropy}[1]{}%
\renewcommand{\minentropy}[1]{\ensuremath{\mathsf{MinEnt}(#1)}}
